\section{A convergent all-orders analytic inverse algorithm}
\label{sec:allinverse}
A further separately audited extension \cite{allinverse,allaudit}
supplies an all-orders analytic inverse with an explicit geometric
tail. It is a Lagrange derivative algorithm in an auxiliary parameter;
it does not claim that every inverse exponential sector has been
rearranged into a canonical ordered transseries.

\subsection{Holomorphic formulas without logarithmic winding}
Fix one branch and a real center $r_0\ge100$. On the closed complex
disk $|r-r_0|\le1/4$, write
\[
 p(r)=\alpha_*(r+\epsilon),\quad L(r)=\lambda_*(r+\epsilon),\quad
 y(r)=L(r)-p(r)-1,
\]
where $(\alpha_*,\lambda_*,\epsilon)$ is
$(2/3,1,1)$, $(2/3,1,-1)$, or $(5/7,5/4,1)$.
Put $p_0=p(r_0)$, $y_0=y(r_0)$, $q_0=2^{-p_0}$.
Here $2^z$ means $\exp(z\ln2)$. Define, in this order,
\begin{align*}
 d_c&=2^{-p}+2\,2^{-p-y},&
 a_c&=2^{-2p-2y}(1+d_c)^{-2},\\
 \ell_A&=L\ln2+\ln(1+d_c)+g(a_c),&
 u_c&=\exp(-2p\ell_A),\\
 b_c&=16u_c^2/(1-u_c)^4,&
 \ell_S&=2p\ell_A-\ln2+2\ln(1-u_c)+g(b_c),\\
 v_c&=\exp(-2r\ell_S).&&
\end{align*}
The logarithms $\ln(1+z)$ and $\ln(1-z)$ are their Taylor branches
at zero, and $g$ is the convergent series from
Theorem~\ref{thm:forward}. These formulas specify holomorphic
logarithms $\ell_A,\ell_S$; they never take a principal
$\log S$ or principal $\operatorname{arcosh}S$ on the complex disk,
where the values of $S$ may wind around zero.
Let $N(r)$ be the branch cubic \eqref{eq:Ncubics}, and define
\begin{align}
 E_c(r)=\frac1{\ln2}\bigl\{&2p(r)(r-1)[\ln(1+d_c)+g(a_c)]
 +2(r-1)\ln(1-u_c)\nonumber\\
 &+r g(b_c)-\tfrac12\ln(1-b_c)+\ln(1-v_c)\bigr\}.
 \label{eq:complexE}
\end{align}
On the real axis this is exactly $\log_2\mathcal H(r)-N(r)$.
We prove all required branch and nonvanishing assertions, rather than
assuming that these formulas are well defined.

\begin{lemma}[Uniform complex-disk bounds]\label{lem:disk}
On a neighborhood of $|r-r_0|\le1/4$, $E_c$ and
\[
 D_c(r;r_0)=\frac{N(r)-N(r_0)}{r-r_0},\qquad
 D_c(r_0;r_0)=N'(r_0),\qquad \mathfrak f(r)=E_c(r)/D_c(r;r_0)
\]
are holomorphic, and on the closed disk
\begin{equation}\label{eq:diskbounds}
 |E_c(r)|<18r_0^2q_0,\qquad |D_c(r;r_0)|>2r_0^2,\qquad
 |\mathfrak f(r)|<\varepsilon_0:=16q_0<1/8.
\end{equation}
\end{lemma}
\begin{proof}
Uniformly, $p_0\ge66$, $y_0\ge32$, and $|p(r)|<r_0$.
The slopes of $p,y$ are positive and less than one, so
\[
 |2^{-p}|\le2^{1/4}q_0<(5/4)q_0,\quad
 |2^{-y}|<(5/4)2^{-y_0},\quad |d_c|<2q_0\le2^{-65}.
\]
Also $|2^{-2p-2y}|<4q_0^2\,2^{-2y_0}$ and
$|1+d_c|^{-2}<2$, whence $|a_c|<q_0^2$.
The convergent series give, with $\Phi=\ln(1+d_c)+g(a_c)$,
\begin{equation}\label{eq:Phi}
 |\ln(1+d_c)|<4q_0,\qquad |g(a_c)|<q_0^2,\qquad
 |\Phi|<5q_0.
\end{equation}

Write $h_c=r-r_0$ and $z_0=r_0+\epsilon$. The inequality
$\operatorname{Re}((z_0+h_c)^2)\ge z_0^2-2z_0|h_c|-|h_c|^2$
and $2/3\le\alpha_*\lambda_*\le25/28$ imply
\[
 \operatorname{Re}(pL)\ge\frac23
 \left(r_0^2-\frac52r_0+\frac7{16}\right)\ge\frac35r_0^2.
\]
The last comparison follows by direct subtraction for $r_0\ge100$.
As $5r_0q_0<1$, \eqref{eq:Phi} gives
\[
 \operatorname{Re}(p\ell_A)\ge\tfrac35r_0^2\ln2-1
 >\tfrac12r_0^2\ln2.
\]
Therefore
\begin{equation}\label{eq:ubdisk}
 |u_c|<2^{-r_0^2}\le q_0^4,\qquad
 |b_c|\le\frac{16q_0^8}{(1-q_0^4)^4}<32q_0^8<q_0^4.
\end{equation}
Here $p_0<r_0$ and $r_0^2\ge4r_0$ justify the first comparison.
This validates the Taylor branches involving $u_c,b_c$. In particular
\[
 \ell_S=2pL\ln2-\ln2+Q_c,\qquad
 Q_c=2p\Phi+2\ln(1-u_c)+g(b_c),\qquad |Q_c|<11r_0q_0.
\]

The cubic $r(r+\epsilon)^2$ differs from its center value by at most
\[
 \tfrac14(3r_0^2+4r_0+1)+\tfrac1{16}(3r_0+2)+\tfrac1{64}<r_0^2.
\]
Since $\alpha_*\lambda_*<1$, it follows that
\[
 \operatorname{Re}(rpL)\ge\tfrac23r_0(r_0-1)^2-r_0^2
 \ge\tfrac35r_0^3.
\]
Using $22r_0^2q_0<1$ and $|r|\le r_0+1/4$ gives
\[
 \operatorname{Re}(r\ell_S)
 \ge\tfrac65r_0^3\ln2-(r_0+1/4)\ln2-1>r_0^3\ln2,
 \qquad |v_c|<2^{-2r_0^3}\le q_0^4.
\]
The elementary inequalities $5r_0q_0<1$ and $22r_0^2q_0<1$
hold at $r_0=100$ using $q_0\le2^{-66}$; their left sides
decrease thereafter because $p_0'\ge2/3$.

Now $|2p(r-1)|<2r_0^2$ and $|r-1|<r_0$.
The first numerator term in \eqref{eq:complexE} is less than
$10r_0^2q_0$. By \eqref{eq:ubdisk}, the remaining terms have
total modulus below
\[
 4r_0q_0^4+2r_0q_0^4+q_0^4+2q_0^4
 =(6r_0+3)q_0^4<r_0^2q_0.
\]
Since $1/\ln2<3/2$, this gives
$|E_c|<(33/2)r_0^2q_0<18r_0^2q_0$.

If $\gamma<2$ is the cubic's leading coefficient, then exactly
\[
 D_c(r_0+h_c;r_0)=N'(r_0)+\tfrac12N''(r_0)h_c+\gamma h_c^2.
\]
The real derivative bounds imply
$|D_c|>3r_0^2-(3/2)r_0-1/8>2r_0^2$.
Thus $D_c$ never vanishes and $|\mathfrak f|<9q_0<16q_0$.
All series arguments have modulus strictly below one, with strict
uniform bounds on the compact disk. Those bounds persist on a
slightly larger disk, proving holomorphy on a neighborhood of its closure.
\end{proof}

\begin{theorem}[Lagrange inverse with an explicit tail]\label{thm:lagrange}
Put $\rho_0=1/(8\varepsilon_0)>1$. The unique small solution of
\begin{equation}\label{eq:himplicit}
 h(z)=-z\,\mathfrak f(r_0+h(z)),\qquad h(0)=0,
\end{equation}
is holomorphic on a neighborhood of $|z|\le\rho_0$ and has
$|h(z)|<1/4$ there. Its convergent Lagrange series is
\begin{equation}\label{eq:lagrange}
 h(z)=\sum_{n\ge1}\frac{(-z)^n}{n!}
 \left.\frac{d^{n-1}}{dr^{n-1}}\bigl[\mathfrak f(r)^n\bigr]
 \right|_{r=r_0}.
\end{equation}
At $z=1$, $r=r_0+h(1)$ is the exact inverse of
$\log_2\mathcal H(r)=N(r_0)$ in the disk. For every $K\ge0$,
\begin{equation}\label{eq:lagrangetail}
 \left|h(1)-\sum_{n=1}^K\frac{(-1)^n}{n!}
 \left.\frac{d^{n-1}}{dr^{n-1}}[\mathfrak f(r)^n]\right|_{r=r_0}\right|
 \le\frac14\frac{\rho_0^{-K-1}}{1-\rho_0^{-1}}.
\end{equation}
\end{theorem}
\begin{proof}
On $|h|=1/4$ and $|z|\le\rho_0$,
$|z\mathfrak f(r_0+h)|<\rho_0\varepsilon_0=1/8<|h|$.
Rouch\'e's theorem gives exactly one zero, counting multiplicity,
of $h+z\mathfrak f(r_0+h)$ in the disk. It is therefore simple.
The analytic implicit-function theorem, together with uniqueness,
patches the local solution into one holomorphic function throughout
the parameter disk and its neighborhood. Lagrange inversion at zero
gives \eqref{eq:lagrange}. At $z=1$, multiplying
\eqref{eq:himplicit} by $D_c(r_0+h;r_0)$ gives
$N(r_0+h)+E_c(r_0+h)=N(r_0)$, the desired inverse equation.
Finally Cauchy's coefficient estimate on $|z|=\rho_0$ is
$|[z^n]h(z)|\le(1/4)\rho_0^{-n}$. Summing its geometric tail
proves \eqref{eq:lagrangetail} and absolute convergence at $z=1$.
\end{proof}

For an actual $R\ge100$, take $N(r_0)=\log_2\mathcal H(R)$
as in Theorem~\ref{thm:inverse}. Its localization gives
$0<r_0-R<2d_0<1/4$, so this real $R$ is the unique root above;
conjugation symmetry also makes $h(1)$ real. At an arbitrary center
$r_0=100$, the inverse of the target $N(100)$ may instead lie
slightly below $100$. The disk theorem then concerns the local
analytic continuation, not an inverse constrained to the half-line
$R\ge100$. This endpoint caveat does not affect the application
starting from an actual $R\ge100$, where $r_0>R\ge100$.

The first Lagrange term is $-E_c(r_0)/N'(r_0)$ and retains all
forward scales. Its leading $2^{-p_0}$ part agrees with
Theorem~\ref{thm:inverse}, while \eqref{eq:lagrangetail} makes
the $n\ge2$ terms uniformly $O(q_0^2)$.
Every finite derivative in \eqref{eq:lagrange} is computable by
power-series arithmetic on the explicitly specified holomorphic
functions. This is a genuinely convergent all-orders algorithm with
a computable error bound. Its ordering is by the artificial parameter
$z$, and it does not smooth the input's radix jumps.

